% modules/kls/reading-paths.tex --- front matter, before Part I.
%
% Unnumbered by design: this is navigation, not mathematics. It carries no claim
% environment and therefore no ledger node, and its filename has no numeric prefix
% because main.tex, not the prefix, fixes reading order.
\documentclass[../../main.tex]{subfiles}
\begin{document}

\section*{How to read this document}
\addcontentsline{toc}{section}{How to read this document}
\label{sec:reading-paths}

This manuscript is long because it is a research record as well as an exposition. It is
meant to be entered at several depths, and only the first of the paths below is meant to
be read in full by everyone.

\begin{description}[leftmargin=0pt,style=nextline]
  \item[A twenty-minute overview.] Part~I, then the atlas in
  Section~\ref{sec:frontier-atlas}. That is the conjecture, the frontier, the obstruction
  that shapes every attempt, and one table per live route giving its advance and its exact
  blocker. Nothing else is needed to say accurately what this project has and has not done.

  \item[The methods of the literature.] Part~II, Sections~\ref{sec:family-needles}
  through~\ref{sec:family-transport}, five families in a common format. Each closes by
  naming the live route it feeds, so the survey can be read either on its own or as the
  approach to Parts~III and~IV.

  \item[The stochastic-localization routes.] The prelude,
  Section~\ref{sec:localization-prelude}, then Route~E from
  Section~\ref{sec:introduction} or Route~S from Section~\ref{sec:spectral-route}. The
  prelude replaces the full apparatus with what a reader needs to follow the argument; the
  apparatus itself is Appendices~\ref{sec:notation}--\ref{sec:models}, linked where used.

  \item[The deterministic moment-map route.] Section~\ref{sec:moment-map-cmh} for the
  programme, Section~\ref{sec:cmh-normalization} for the endpoint and the proof that it
  dominates the affine Poincar\'e constant, Section~\ref{sec:cmh-exact-cases} for the cases
  where it is computed exactly. No stochastic localization is used anywhere in it.

  \item[The open bottlenecks, as a list.] Section~\ref{sec:open} for Route~E,
  Section~\ref{subsec:spectral-headline} for Route~S, and
  Section~\ref{sec:kls-synthesis} for the four cross-route targets, each mapped to the
  ledger node that carries it.

  \item[The complete technical development.] Front to back, with the appendices read at
  the point they are first linked rather than in advance.
\end{description}

\paragraph{Reading the status labels.}
Every theorem, lemma, question and obstruction below carries a small standing label next to
its statement --- \emph{published result}, \emph{unreviewed preprint}, \emph{certified
here}, \emph{open question}, and so on. These are generated from
\texttt{research/program/ledger.yaml} and are not maintained by hand, so a statement's
typography never outruns its evidence. Two distinctions they preserve are easy to lose and
matter throughout: an unreviewed preprint is recorded as \emph{open} however plausible it
is, and a \emph{proved} implication whose antecedent is still open is proved but not
progress on the conjecture.

\end{document}
